\chapter{Results}
\label{ch:results}

The delivered forecaster reaches a mean absolute error of 8,796~kWh per half-hour over the
test period, 37.4\% below \naiveA{} and 56.3\% below \naiveB{}.
Table~\ref{tab:headline} gives that result in full.

\begin{table}[htp]
\centering
\begin{tabular}{lrrrrr}
\toprule
& {\bf MAE} & {\bf RMSE} & {\bf MASE} & {\bf sMAPE} & {\bf Coherence} \\
\midrule
\naiveA{} & 14,053 & 28,273 & 1.149 & 32.3 & 0 \\
\naiveB{} & 20,117 & 38,419 & 2.011 & 41.9 & 0 \\
Delivered model, unreconciled & 9,240 & 18,593 & 0.796 & 23.9 & 107,623 \\
Delivered model with MinT & {\bf 8,796} & 16,954 & 0.821 & 24.7 & {\bf 0} \\
\bottomrule
\end{tabular}
\caption[Headline test-period result]{\textbf{Test-period performance of the delivered
forecaster.} MAE and RMSE are in kWh per half-hour, pooled over all
$19 \times 107 \times 48 = 97{,}584$ forecast values. Coherence is the largest absolute
discrepancy between a parent node's forecast and the sum of its children, so zero means
the hierarchy is exactly self-consistent.}
\label{tab:headline}
\end{table}

That single number is the end of a chain, and the chain is what this chapter reports. The
experiments of Chapter~\ref{ch:experiments} were not a survey of independent questions run
in parallel; each one was prompted by the result of the one before it, and several were
prompted by a result that contradicted what had just been concluded. The sections below
follow that order, and each closes with the question it hands to the next.

The chain runs as follows. A first attempt at covariate selection gave an answer that the
test period reversed, which forced the protocol to be rebuilt before any substantive
question could be asked (Section~\ref{sec:res-protocol}). The rebuilt protocol then had to
be calibrated, since a decision rule needs to know what counts as a difference. With a
protocol in hand, the obvious first question was whether the forecasting problem wants a
sophisticated architecture, and the answer from nine deep candidates was that it does not
(Section~\ref{sec:res-arch}); that section also produced two results which changed how
everything afterwards had to be measured and reported. Attention then moved to what the
model is given rather than what the model is, and the covariate experiments found that
restricting the information helped while accumulating it did not
(Section~\ref{sec:res-cov}). Those experiments left the hierarchy internally inconsistent,
which reconciliation was introduced to repair, at a price that only two of the four
projections earn back (Section~\ref{sec:res-recon}). Two questions raised by the
supervisors then had to be answered on the finished pipeline: whether it extends to a
fortnight (Section~\ref{sec:res-horizon}), and what had been left on the table by never
tuning the learner (Section~\ref{sec:res-tuning}). The last of those uncovered a defect in
the training code which turned out to be worth more than the tuning it was found during.

%---------------------------------------------------------------------------------------
\section{The protocol, and the scale it established}
\label{sec:res-protocol}

\subsection{Why the protocol had to come first}

The first version of this work ranked eleven candidate feature sets on a single contiguous
validation block. Re-ranking the same eleven on the test period reversed the ordering
completely: the family that proved most valuable, per-station weather at the generating
plants, had ranked eighth of eleven on validation and ranked first on test. The cause was
seasonal rather than statistical. The validation block spanned summer, and almost all of
the value of weather information sits in winter, when wind and inflow drive dispatch. One
block measures the season it happens to cover.

A result that inverts under a change of evaluation window is not a result, so nothing else
could be concluded until the evaluation itself was fixed. Four-fold cross-season validation
replaced the single block, and every figure in this chapter rests on it.

\subsection{What counts as a difference}

A four-fold protocol still cannot decide anything without a scale. Re-running identical
configurations across the folds, and re-running configurations whose only difference was
execution order, put that scale at roughly $\pm 0.6$~\pp{} on the 19-node median.
Differences smaller than that carry no information about the variable under test, however
suggestive their sign.

That number is why the decision rule rests on sign agreement across folds rather than on
the four-fold mean. A mean of $-0.3$~\pp{} says nothing on its own. Four folds that
independently agree in sign say a great deal, since under pure noise that happens one time
in sixteen, and the strength of that argument does not depend on the size of the effect.
Several of the changes adopted later in this chapter have means inside the noise floor and
were accepted on sign agreement, then confirmed on the test period. Each case is flagged
where it arises.

With an evaluation that survives a change of window and a rule for reading it, the first
substantive question could be asked: does this problem call for a sophisticated model?

%---------------------------------------------------------------------------------------
\section{Architecture: the answer is no, and two complications}
\label{sec:res-arch}

\subsection{None of the nine candidates beat the naive benchmark}

Table~\ref{tab:arch} gives each of the nine trained architectures on its own. The spread
is wide, from a linear model that matches \naiveA{} to a patched transformer three times
worse than it, but the ordering matters less than the fact that no architecture clears the
benchmark. NLinear, the simplest candidate, comes closest.

\begin{table}[htp]
\centering
\begin{tabular}{lrrrr}
\toprule
& \multicolumn{2}{c}{\bf No covariates} & \multicolumn{2}{c}{\bf With 24 covariates} \\
\cmidrule(lr){2-3}\cmidrule(lr){4-5}
{\bf Architecture} & {\bf MAE} & {\bf vs \naiveA{}} & {\bf MAE} & {\bf vs \naiveA{}} \\
\midrule
NLinear & 13,990 & $-0.5$\% & 14,640 & $+4.2$\% \\
TCN & 13,994 & $-0.4$\% & 18,311 & $+30.3$\% \\
RNN & 14,115 & $+0.4$\% & 16,413 & $+16.8$\% \\
LSTM & 14,723 & $+4.8$\% & 24,410 & $+73.7$\% \\
GRU & 15,510 & $+10.4$\% & 19,861 & $+41.3$\% \\
TiDE & 17,021 & $+21.1$\% & 56,067 & $+299.0$\% \\
THP-Former & 18,503 & $+31.7$\% & 21,244 & $+51.2$\% \\
TSMixer & 19,549 & $+39.1$\% & 19,969 & $+42.1$\% \\
PatchTST & 30,511 & $+117.1$\% & 25,079 & $+78.5$\% \\
\midrule
\naiveA{} & 14,053 & --- & & \\
Per-node selection & {\bf 12,873} & {\bf $-8.4$\%} & & \\
\bottomrule
\end{tabular}
\caption[Nine architectures]{\textbf{The nine candidate architectures, each on its own.}
A positive percentage means the architecture is worse than \naiveA{}. Candidates are
trained inside the expert system, where early stopping ends three months earlier than in a
standalone run and no reconciliation is applied, so these figures understate each
architecture slightly: THP-Former trained standalone to completion reaches 14,986.}
\label{tab:arch}
\end{table}

Two things follow immediately. The heavy use of \naiveA{} by the selector is not a defect
of the selector: six of the nineteen nodes choose the untrained fallback because the
candidate pool holds nothing better for those nodes. And per-node selection is where this
system's accuracy comes from, 12,873 against 13,990 for the best single architecture.
Choosing a different model for each node is worth roughly eight percentage points, far more
than the choice among the nine.

The answer to the question that opened this section is therefore no. Capacity is not the
binding constraint on this problem, and the work that follows spends its effort on what the
model is given rather than on what the model is.

Before leaving the architecture experiments, two of their results have to be settled,
because both of them govern how every later result is measured and reported.

\subsection{Complication one: the advantage depends on which aggregation is quoted}

\begin{table}[htp]
\centering
\begin{tabular}{lrr}
\toprule
& {\bf Pooled MAE} & {\bf MASE} \\
\midrule
\naiveA{} & 14,053 & 1.149 \\
Per-node selection & 12,873 & 1.149 \\
With learned BUN & 13,102 & 1.203 \\
\bottomrule
\end{tabular}
\caption[Two aggregations of the expert system]{\textbf{The same forecasts under two
aggregations.} Pooled MAE weights each forecast value equally, so the three large aggregate
nodes dominate. MASE normalises each node by its own training-period naive error and then
weights the nodes equally. The 8.4\% advantage visible in the first column is absent from
the second.}
\label{tab:es-agg}
\end{table}

A MASE of 1.149 against the benchmark's 1.149 is not a near-tie but an exact one. All of
the pooled-MAE advantage sits in the three large aggregate nodes, and once each node is
normalised by its own difficulty and the nodes are weighted equally, nothing remains.

This is the empirical reason for the rule stated in Section~\ref{sec:protocol} that the two
aggregations must never be mixed. A reader given only the first column would conclude that
per-node selection beats the benchmark by 8.4\%; a reader given only the second would
conclude that it exactly matches it. Both readings are arithmetically correct. Every
percentage quoted from here on therefore names its aggregation, and where the two disagree,
both are given.

The learned bottom-up network in the third row removes the hierarchical inconsistency
completely, bringing a maximum parent--child discrepancy of 1,258,079 down to zero, and
raises pooled MAE by 1.8\% and MASE from 1.149 to 1.203 in doing so. That reconciliation
can cost accuracy rather than improve it is the first appearance of a trade-off taken up in
full in Section~\ref{sec:res-recon}.

\subsection{Complication two: informative covariates can still do damage}

The second complication is sharper, because it contradicts what
Section~\ref{sec:res-cov} goes on to establish. Supplying the same 24 covariate columns
that help the gradient-boosting pipeline makes this system worse. Eight of the nine
architectures degrade, several of them severely, and the assembled system falls from
$-8.4$\% to $+6.4$\% before reconciliation.

Two explanations were available: the columns carry no usable information, or the columns are
informative and something about the path they travel into this particular system spoils
them. The diagnostic runs were built to separate the two, by supplying covariates in order
of how much of them is missing.

\begin{table}[htp]
\centering
\begin{tabular}{llrr}
\toprule
{\bf Covariates supplied} & {\bf Missingness} & {\bf Unreconciled} & {\bf Reconciled} \\
\midrule
None & --- & $-8.4$\% & $-6.8$\% \\
Plant availability only (1 column) & none at all & $-5.9$\% & $-2.9$\% \\
Plant weather (10 columns) & 9 of 19 nodes empty & $+14.5$\% & $+0.5$\% \\
All (24 columns) & extensive & $+6.4$\% & $+2.8$\% \\
\bottomrule
\end{tabular}
\caption[Diagnostic ladder]{\textbf{The diagnostic ladder, ordered by how much is
missing.} Percentages are relative to \naiveA{}, and positive is worse than the benchmark.
The second row supplies a single column that is complete for all 19 nodes and requires no
imputation whatsoever.}
\label{tab:es-diag}
\end{table}

The second row decides it. One column, complete everywhere, nothing imputed, and the system
still degrades. A paired comparison against the no-covariate run on the same 107 origins
puts the damage at $+2.7$\% [$+0.3$, $+5.6$] unreconciled and $+4.1$\% [$+1.9$, $+6.6$]
after reconciliation, and neither interval contains zero.

Imputation of missing values is therefore excluded. The explanation lies in how external
information enters this system rather than in the information itself, since the same columns
are measurably useful in Section~\ref{sec:res-cov}. Identifying the responsible property
lay outside the scope of the experiments reported here, and the candidate explanations are
weighed in Chapter~\ref{ch:discussion}.

The consequence for the chain is narrow but real. A covariate experiment can fail for
reasons that have nothing to do with the covariate, so a negative covariate result is only
interpretable within one system at a time. The covariate work of the next section is
accordingly conducted entirely within the gradient-boosting pipeline, and its conclusions
are claimed for that pipeline rather than for the covariates in general.

%---------------------------------------------------------------------------------------
\section{Covariates: restriction beats accumulation}
\label{sec:res-cov}

Since model capacity is not the constraint, the question becomes which information to
supply and where to attach it. The natural starting point is to supply everything
available to every node and let the learner sort it out. That turns out to be the worst
configuration tested.

\subsection{The selection ladder}

Round~19 rebuilt the selection process as a ladder, each rung adding one stage, all five
rungs executed inside a single run so that they share a baseline.

\begin{table}[htp]
\centering
\begin{tabular}{lrrrrrc}
\toprule
{\bf Configuration} & {\bf Feat.} & {\bf F1} & {\bf F2} & {\bf F3} & {\bf F4} & {\bf Folds neg.} \\
\midrule
No covariates & 19 & --- & --- & --- & --- & baseline \\
All weather and hydrology, shared & 63 & $-4.95$ & $-3.51$ & $+1.40$ & $+0.58$ & 2 of 4 \\
Specified base plus wind vectors & 49 & $-4.72$ & $-3.24$ & $-2.20$ & $-1.81$ & 4 of 4 \\
\quad with per-type gating & 49 & $-4.57$ & $-3.90$ & $-1.72$ & $-1.82$ & 4 of 4 \\
\quad plus aggregate-node weather & 67 & $-4.59$ & $-4.45$ & $-5.25$ & $-1.84$ & 4 of 4 \\
\quad plus Meridian hydro (delivered) & 84 & $-6.02$ & $-5.82$ & $-4.73$ & $-2.69$ & 4 of 4 \\
\bottomrule
\end{tabular}
\caption[Selection ladder]{\textbf{Each stage of covariate selection, relative to the
no-covariate baseline.} Figures are \pp{} of skill against \naiveA{}, and negative is
better. Each fold's figure is the median over 19 nodes of the per-node percentage, so nodes
are weighted equally. The second column is the number of features the configuration
carries.}
\label{tab:ladder}
\end{table}

Giving every node every available covariate is the only configuration on the ladder that
fails the sign test, with two folds improving and two deteriorating. Cutting the pool back
to the supervisor-specified base, with wind decomposed into vector components, drops
fourteen features and converts that inconsistent result into a consistent three percentage
points in every fold. More information made the model worse and discipline about which
information reaches which node made it better, which sets the direction for the rest of the
section: the remaining gains come from placing features correctly and from removing families
that do not earn their place, not from finding more data.

The third rung raises a problem rather than settling one. Per-type gating changes the
19-node median by 0.01~\pp{}, and comparing the two gating modes against each other gives
the same verdict, $-25.30$\% under \verb+strict+ against $-25.22$\% under \verb+dispatch+.
On this metric, gating does nothing. Since gating was one of the two constraints imposed at
the outset, a result of nothing needs explaining before the ladder can be trusted, and the
explanation turns out to be the metric rather than the gating.

\subsection{Why the 19-node median is the wrong instrument}
\label{sec:res-cov-cat}

The supervisors required that results be read per generator type rather than as a single
median. Splitting the first elimination pass that way shows the requirement doing real
work.

\begin{table}[htp]
\centering
\begin{tabular}{lrrrrr}
\toprule
{\bf Family removed} & {\bf Wind (2)} & {\bf Hydro (9)} & {\bf Geo. (6)} & {\bf Top agg. (2)} & {\bf Median} \\
\midrule
Aggregate-node weather & $+0.37$ & $+1.26$ & $+2.34$ & $+5.02$ & $+1.14$ \\
POCP, three columns & $-0.31$ & $+0.89$ & $+4.26$ & $+1.55$ & $+1.30$ \\
Earlier storage version & $+0.34$ & $-3.35$ & $+3.65$ & $-1.89$ & $-1.43$ \\
Registered capacity & $+0.42$ & $+0.75$ & $-0.76$ & $+0.49$ & $+0.49$ \\
National demand & $-0.87$ & $+0.55$ & $-0.72$ & $-0.40$ & $-0.02$ \\
\bottomrule
\end{tabular}
\caption[Elimination by generator category]{\textbf{The first elimination pass split by
generator type.} Node counts are given in the header, and a positive figure means removal
hurt that group.}
\label{tab:bycat}
\end{table}

Removing the three POCP columns costs the six geothermal nodes 4.26\%, the largest single
effect anywhere in the table, while the median records 1.30\% with an inconsistent sign.
Maintenance scheduling matters to geothermal plant, which runs at near-constant output until
it is taken offline, and barely matters to hydro, whose output varies for other reasons.
Averaged over nineteen nodes, thirteen of which the family does not concern, a strong effect
on six nodes becomes a weak and unreliable one.

The earlier storage family shows the same arithmetic with opposite signs. Removing it helps
hydro by 3.35\% and harms geothermal by 3.65\%; the two cancel in the median to $-1.43$\%
and conceal that the family is actively mis-specified for one of the two groups.

So the constraint handed down at the outset was not a stylistic preference. A single median
over a hierarchy whose leaves run on different physics will hide effects of several per
cent, and gating registering zero on the median says nothing about whether gating works.
It also sets up the next step: if features have to be placed per node type, then the four
nodes that have no physical plant of their own need features of their own.

\subsection{Aggregate nodes need features built for them}

Weather is read at plant coordinates, which the four aggregate nodes do not have. The
\verb+agg_+ family supplies them with capacity-weighted averages of their leaf descendants'
weather, with wind converted to vector components before averaging and the capacity weights
recomputed at every timestamp.

Removing that family costs the four aggregate nodes 6.46, 6.79, 3.34 and 2.11~\pp{} of
skill in the four folds, all in the same direction, pooling to 4.18. The fifteen leaf nodes
are unaffected, which is what the construction predicts: the family reaches only the
aggregate nodes, so it should move only those. Carried to the test period, the same
configuration loses 5.17\% of pooled MAE with a paired interval of [$+3.39$, $+6.83$], and
each of the three principal aggregate nodes gives up more than 9\%.

This is the one covariate result that is large, consistent across folds, confirmed out of
sample, and mechanically explicable at the same time, and it is also the result that most
directly rewards taking the hierarchy seriously as structure rather than as nineteen
separate series.

\subsection{Five of six families do not earn their place}

If placing features well is what pays, the obvious test is how many of the families being
carried are doing any work at all. Two elimination passes removed each family from the full
set in turn and measured the damage.

\begin{table}[htp]
\centering
\begin{tabular}{lrrrrrc}
\toprule
{\bf Family removed} & {\bf F1} & {\bf F2} & {\bf F3} & {\bf F4} & {\bf Mean} & {\bf Verdict} \\
\midrule
\multicolumn{7}{l}{\emph{First pass}} \\
Aggregate-node weather & $+1.49$ & $+1.63$ & $+0.61$ & $+0.82$ & $+1.14$ & keep \\
POCP, three columns & $+0.34$ & $+1.58$ & $+3.66$ & $-0.40$ & $+1.30$ & inconsistent \\
Registered capacity & $+0.68$ & $+0.79$ & $+0.51$ & $-0.03$ & $+0.49$ & inconsistent \\
National demand & $+0.35$ & $+0.36$ & $+0.08$ & $-0.87$ & $-0.02$ & inconsistent \\
Earlier storage version & $+0.52$ & $+0.95$ & $-6.17$ & $-1.01$ & $-1.43$ & inconsistent \\
\midrule
\multicolumn{7}{l}{\emph{Second pass, on the survivors}} \\
Aggregate-node weather & $+0.85$ & $+0.14$ & $+0.26$ & $+1.23$ & $+0.62$ & keep \\
POCP & $+0.26$ & $-0.49$ & $-0.35$ & $+0.25$ & $-0.08$ & inconsistent \\
Registered capacity & $-0.15$ & $-0.48$ & $+0.16$ & $+0.44$ & $-0.01$ & inconsistent \\
National demand & $-0.02$ & $+0.23$ & $+0.32$ & $+0.11$ & $+0.16$ & inconsistent \\
\bottomrule
\end{tabular}
\caption[Backward elimination]{\textbf{Two passes of backward elimination.} A positive
figure means removal hurt, so the family was doing work. Only aggregate-node weather hurts
in all four folds, in both passes.}
\label{tab:backward}
\end{table}

Exactly one family of six clears the sign rule, in both passes. Several of those that fail
it have large four-fold means, and the POCP row of the first pass has a larger mean than the
family that was kept, yet one fold moves the other way. The previous subsection explains
why that is the right verdict and not merely a cautious one: POCP matters intensely to six
nodes and the median cannot see it, so the median's inconsistency is real information about
the median, not about POCP.

\subsection{Testing the principle by adding instead of removing}
\label{sec:res-cov-neg}

Everything so far points one way, that this pipeline wants few and well-placed features.
A principle established by removing things should be tested by adding things, and
Rounds~21 and~22 did that. The asymmetry noted in Section~\ref{sec:exp-negative} governs
how their results read: for an addition, insufficient evidence means do not add.

Round~21 tested cloud cover, radiation, dew point, evapotranspiration and soil moisture in
three combinations, with four-fold means of $+0.09$, $+0.31$ and $+0.36$~\pp{}. All three
have inconsistent signs, all three are mildly harmful on average, and all three sit inside
the $\pm 0.6$~\pp{} noise floor of Section~\ref{sec:res-protocol}. Insufficient evidence,
so the five variables stay out.

Round~22 produced a stronger verdict, and a more instructive one.

\begin{table}[htp]
\centering
\begin{tabular}{lrrrrrc}
\toprule
{\bf Group} & {\bf F1} & {\bf F2} & {\bf F3} & {\bf F4} & {\bf Mean} & {\bf Verdict} \\
\midrule
19-node median & $+0.39$ & $+0.71$ & $+2.83$ & $+0.48$ & $+1.10$ & eliminate \\
Hydro (9) & $+1.22$ & $+1.51$ & $+3.64$ & $+0.75$ & $+1.78$ & harmful \\
Geothermal (6) & $-3.11$ & $-1.60$ & $-1.62$ & $-0.49$ & $-1.71$ & helpful \\
Wind (2) & $+0.62$ & $+0.66$ & $+0.74$ & $-1.31$ & $+0.18$ & inconsistent \\
Top aggregates (2) & $-1.28$ & $+0.33$ & $+3.48$ & $+0.72$ & $+0.81$ & inconsistent \\
\bottomrule
\end{tabular}
\caption[Snowpack by category]{\textbf{South Island snowpack, by generator category.}
Positive means adding the family hurt. The family was motivated by an effect on hydro
scheduling, and hydro is the group it damages most.}
\label{tab:snow}
\end{table}

The snowpack family rests on a specific physical argument. South Island snowmelt feeds South
Island inflow, which affects national dispatch, which affects where Mercury's hydro stations
sit in the merit order. If that chain operated, hydro is where the benefit should appear.
Hydro is instead the worst-affected group, harmed in all four folds and by 1.78\% on
average, so the motivation is not borne out. The family does help the six geothermal nodes
in all four folds, which was not predicted and which no mechanism available to the
experiment supports; geothermal output is largely independent of hydrology, so the most
economical reading is that three slowly varying columns act as a weak seasonal proxy for
those nodes, a role already filled by features the model carries. On that reading the result
belongs to this dataset rather than to the physics, and it was not pursued.

Building the family was prompted by an observation that explains why an existing family had
been weaker than expected: the original snow features covered the Taupo and Waikato
catchments, which have had essentially no snow since 2021, so for the whole study period
that family had been supplying a constant plus a slow trend rather than snow information.

The pair of rounds completes the covariate argument in the form it needs. Restriction
helped, placement helped, and the two attempts to enlarge the pool failed, one for want of
evidence and one against its own stated mechanism. What the covariate work leaves behind is
a model that is accurate node by node and inconsistent as a hierarchy, with a maximum
parent--child discrepancy of 107,623~kWh. Repairing that is the next link.

%---------------------------------------------------------------------------------------
\section{Reconciliation: coherence has a price}
\label{sec:res-recon}

\subsection{Two of four projections are worse than doing nothing}

\begin{table}[htp]
\centering
\begin{tabular}{lrrrr}
\toprule
{\bf Method} & {\bf MAE} & {\bf vs \naiveA{}} & {\bf Coherence} & {\bf Negatives} \\
\midrule
Unreconciled & 9,240 & $-34.2$\% [$-37.7$, $-30.5$] & 107,623 & 0\% \\
Bottom-up & 9,525 & $-32.2$\% [$-36.0$, $-28.3$] & 0 & 0\% \\
OLS & 9,601 & $-31.7$\% [$-35.4$, $-27.6$] & 0 & 0.970\% \\
WLS & 8,922 & $-36.5$\% [$-40.0$, $-32.9$] & 0 & 0.400\% \\
MinT & {\bf 8,796} & {\bf $-37.4$\% [$-40.5$, $-34.1$]} & 0 & 0.870\% \\
\bottomrule
\end{tabular}
\caption[Four reconcilers]{\textbf{The four reconcilers on the delivered base forecasts.}
Intervals are paired bootstrap over the 107 forecast origins. The final column is the
fraction of reconciled values below zero.}
\label{tab:recon}
\end{table}

Bottom-up and OLS both end up worse than the forecasts they were applied to. The coherence
constraint moves error from nodes that were estimated well onto nodes that were estimated
badly, and a projection that ignores how reliable each node is pays that cost without
getting anything back. Weighting by residual variance recovers it and weighting by the full
residual covariance recovers slightly more. The ordering held every time it was re-checked
after a change of base model, which matters because the covariance is re-estimated on each
occasion and the ordering was not guaranteed to survive.

The first appearance of this trade-off was in Section~\ref{sec:res-arch}, where a learned
bottom-up network bought coherence at a cost of 1.8\% of MAE. Table~\ref{tab:recon} shows
that the cost is not inherent to reconciliation but specific to projections that treat the
nodes as equally reliable.

\subsection{Where the price is paid}

\begin{table}[htp]
\centering
\begin{tabular}{lrrrrr}
\toprule
{\bf Level} & {\bf Unrec.} & {\bf BU} & {\bf OLS} & {\bf WLS} & {\bf MinT} \\
\midrule
Company & 0.835 & 0.991 & 0.841 & 0.861 & {\bf 0.838} \\
Island & 0.636 & 0.696 & 0.717 & 0.634 & {\bf 0.633} \\
Fuel type & 0.901 & 0.772 & 1.086 & 0.773 & {\bf 0.761} \\
Station & {\bf 0.807} & {\bf 0.807} & 1.051 & 0.844 & 0.856 \\
\bottomrule
\end{tabular}
\caption[MASE by level]{\textbf{MASE by hierarchy level.} MinT is best at the three upper
levels and worse than no reconciliation at all at the station level.}
\label{tab:bylevel}
\end{table}

MinT wins at the company, island and fuel-type levels and loses at the station level, where
leaving the forecasts alone is better. The burden has to land somewhere, and the leaves are
where it has to land, since the aggregates are sums of the leaves. A trading application
concerned with aggregate volumes should use MinT; one concerned with individual units would
reach the opposite conclusion from the same table. This is the same lesson as
Section~\ref{sec:res-arch}'s two aggregations, arriving from a different direction: the
question ``which method is best'' is incomplete until the level and the metric are named.

\subsection{An open defect, and an attempt to improve the weights}

Every projection that weights nodes unequally produces negative generation: 0.970\% of
values under OLS, 0.870\% under MinT, 0.400\% under WLS. The base forecasts contain none,
because they are clipped at zero during inference, and bottom-up contains none, because it
only adds non-negative leaf values together. Negative output is physically impossible, so
this is a defect of the delivered system rather than a curiosity.

Nothing attempted in this work moved it. The hyperparameter search of
Section~\ref{sec:res-tuning} left it at 0.870\% against 0.889\% before tuning, and the
early-stopping repair in that section, which produced a genuinely out-of-sample residual set
for the covariance estimate and might have been expected to help, left it there too. All
three weighted projections are unconstrained linear maps, and an unconstrained linear map of
non-negative inputs can return negative outputs. Removing the negatives needs reconciliation
under an explicit non-negativity constraint, a quadratic program rather than a closed form,
which was not attempted here.

A separate attempt to improve the weights themselves started from the observation that a
single covariance matrix spans all 48 forecast steps although the error structure at one
step ahead and at 48 steps ahead is not the same.

\begin{table}[htp]
\centering
\begin{tabular}{lrr}
\toprule
{\bf Grouping} & {\bf vs pooled} & {\bf Negatives} \\
\midrule
Pooled (current) & --- & 0.920\% \\
Blocks of 16 steps & $-0.52$\% [$-1.21$, $+0.16$] & 0.808\% \\
Blocks of 8 steps & $-0.56$\% [$-1.47$, $+0.29$] & 0.756\% \\
Blocks of 4 steps & $-0.87$\% [$-2.18$, $+0.36$] & 0.723\% \\
One per step & $-0.46$\% [$-2.68$, $+1.65$] & 0.609\% \\
\bottomrule
\end{tabular}
\caption[Horizon-grouped weights]{\textbf{Estimating the reconciliation weights separately
by forecast horizon.} Every point estimate favours grouping and no interval excludes zero.}
\label{tab:horizgroup}
\end{table}

All four granularities improve on the pooled matrix, the improvement holds for both base
models tested, and the eight resulting cells agree in sign. No interval excludes zero and
the per-origin win rate never exceeds 58\%, so the verdict is insufficient evidence and the
grouping was not adopted. Applying the same grouping to the WLS weights instead changes
nothing at all, by at most 0.1\% and with inconsistent sign, which locates whatever effect
exists in the correlation structure between nodes rather than in the individual variances.

The granularity was selected inside the validation period, split chronologically in half,
fitting on one half and scoring on the other in both directions, and that procedure chose
one matrix per step. On the test period the finest granularity is not the best; blocks of
four steps are. Both halves come from the same three-month window, so a per-step
idiosyncrasy of that window transfers between the halves while failing to transfer to a
later window, and a split that shares a window cannot detect it. Choosing a granularity
properly needs a third independent window, which this design does not have, and using the
test period for the choice would have been selection on the test set. The defect is in the
selection design and it is reported here rather than in the discussion because it bounds how
far Table~\ref{tab:horizgroup} can be trusted.

With a coherent hierarchy and a reconciler chosen, the pipeline was complete enough to
answer the two questions the supervisors had asked of it.

%---------------------------------------------------------------------------------------
\section{Horizon: the long-range gain is not available in practice}
\label{sec:res-horizon}

The first question was whether the pipeline extends from the next day to the next
fortnight. It does, but the covariates stop paying for themselves, and whether they pay at
all depends on an assumption about what a forecaster knows.

\begin{table}[htp]
\centering
\begin{tabular}{lrr}
\toprule
{\bf Horizon} & {\bf Read at origin only} & {\bf Read at target time} \\
\midrule
24 hours (48 steps) & --- & current configuration \\
Two weeks (672 steps) & $+1.0$\% & $-3.5$\% \\
\bottomrule
\end{tabular}
\caption[Horizon extension]{\textbf{Extending the horizon to two weeks.} Figures are the
change in MAE from adding the covariates at that horizon. Negative is an improvement, so
$+1.0$\% means the covariates are a net burden.}
\label{tab:horizon}
\end{table}

Read only at the forecast origin, which is the discipline a live system must obey, the
covariates become a net burden at two weeks. Read at the target time, which requires weather
information that a forecaster standing at the origin does not possess, they help by 3.5\%.
The entire covariate benefit at long range therefore rests on information that is not
available when the forecast is issued. At 24 hours the gap between the two readings is small
enough not to affect any conclusion in this chapter, and the supervisors agreed on
2026-09-11 that forecast vintage would be set aside at this stage.

%---------------------------------------------------------------------------------------
\section{Hyperparameters, and a defect found while tuning}
\label{sec:res-tuning}

The second question was what had been left on the table by never tuning the learner. The
answer is a modest amount, and the search turned up something larger on the way.

\subsection{Round one: two parameters, and a trap in the evidence}

The hand-set configuration ranked fifth of the sixteen sampled, and the three that beat it
agreed on two settings: each used the lowest learning rate on the grid and the largest
minimum leaf size. Both point the same way, that the original configuration trained too
aggressively, and the selected number of boosting rounds agrees, rising from 3,073 for the
incumbent to between 4,900 and 6,000 for the leaders.

\begin{table}[htp]
\centering
\begin{tabular}{lrrrrrc}
\toprule
{\bf Configuration} & {\bf F1} & {\bf F2} & {\bf F3} & {\bf F4} & {\bf Mean} & {\bf Folds neg.} \\
\midrule
$\eta$ 0.03, leaf 400, feat.\ 0.70 & $-0.80$ & $-0.50$ & $-0.17$ & $-0.43$ & $-0.47$ & 4 of 4 \\
$\eta$ 0.03, leaf 400, feat.\ 1.00 & $-0.47$ & $-0.38$ & $-0.20$ & $+0.29$ & $-0.19$ & 3 of 4 \\
$\eta$ 0.03, 63 leaves, leaf 400 & $-0.47$ & $+0.43$ & $-0.35$ & $-0.13$ & $-0.13$ & 3 of 4 \\
Hand-set incumbent & $0.00$ & $0.00$ & $0.00$ & $0.00$ & $0.00$ & 0 of 4 \\
\bottomrule
\end{tabular}
\caption[Round-one confirmation]{\textbf{Round one on all four folds,} as \pp{} of skill
relative to the hand-set configuration. Only the first configuration improves in every
fold; the other two have better-looking means than their fold pattern justifies.}
\label{tab:tune1}
\end{table}

The winner's four-fold mean of $-0.47$~\pp{} is smaller than the noise floor established in
Section~\ref{sec:res-protocol}, which needs stating plainly rather than being left for a
reader to notice. The noise floor governs whether a single reading can be interpreted; the
decision here rests on four independent folds agreeing in sign, a different argument whose
strength does not depend on the size of the effect, and it was confirmed afterwards on the
test period with an interval that excludes zero. The effect is small and it is real.

Only two of the six parameters are supported. The other four moved with the winning
configuration without evidence behind them, and the marginal means that appear to support
one of them do not survive a check: the apparent 0.71~\pp{} effect of the bagging fraction
collapses to 0.13~\pp{} once the learning rate is held fixed, because the random sample had
paired the two. The lesson that marginal means from a small random sample are uninterpretable
until co-occurrence has been checked is taken up in Chapter~\ref{ch:discussion}.

\subsection{Round two: one plateau and one gradient}

Both supported parameters sat at the edge of the round-one grid, which usually means the
optimum lies outside it. Extending the grid showed the two directions behaving quite
differently.

The minimum leaf size and the leaf count are flat. At a learning rate of 0.03, four
configurations spanning leaf sizes from 400 to 1,600 and leaf counts of 127 and 255 fall
within a range of 0.30~\pp{} and in no particular order, so the round-one winner sat at the
edge of a plateau rather than on a gradient.

The learning rate is a gradient and it continues past the round-one grid. Lowering it from
0.03 to 0.005 gains a further $-0.50$~\pp{} across all four folds, at the cost of 25,000
boosting rounds instead of 3,700. Two cells in that confirmation hit the 40,000-round
ceiling without early stopping triggering, both at the lowest learning rate on the largest
fold; truncation leaves a model undertrained and so understates its score, which means the
winner's four-fold agreement stands and is conservative while the runner-up's single
positive fold cannot be relied upon.

\subsection{The round-two winner failed on deployment, and why}

Trained as the delivered model and scored on the test period, the round-two winner was
worse than the round-one winner rather than better, on the validation period as well as on
the test period. Either the cross-validation result was wrong or something differed between
measuring a configuration and deploying it.

The quantity that settles this is the product of the selected number of boosting rounds and
the learning rate, which measures how far a model actually trained. Under cross-validation
the two learning rates receive comparable budgets, 110.8 and 126.8, so the comparison is
between learning rates. Under the delivery script the budgets are 45.9 and 30.8, so the
lower learning rate receives a third less training than the higher one and the comparison is
confounded with how far each model trained. The reason is that the two scripts stop on
different data: the delivery script on a separate validation block of 105 origins, the
tuning scripts on the last 15\% of the training window, roughly 970 origins. A small
stopping set gives a noisy validation curve that triggers earlier, and the effect grows as
the learning rate falls. The deployed model was undertrained rather than outperformed.

Aligning the two rules improved both configurations, and improved the starved one by about
twice as much.

\begin{table}[htp]
\centering
\begin{tabular}{llrr}
\toprule
{\bf Configuration} & & {\bf Unreconciled} & {\bf With MinT} \\
\midrule
$\eta = 0.03$ & aligned vs original & $-1.03$\% [$-1.66$, $-0.42$] & $-1.22$\% [$-1.90$, $-0.56$] \\
$\eta = 0.005$ & aligned vs original & $-2.48$\% [$-3.52$, $-1.43$] & $-2.05$\% [$-2.88$, $-1.22$] \\
\bottomrule
\end{tabular}
\caption[Effect of aligning early stopping]{\textbf{Aligning the delivery script's early
stopping with the cross-validation rule.} All four intervals exclude zero. The lower
learning rate, which the original rule starved more severely, gains roughly twice as much.}
\label{tab:esfix}
\end{table}

That the gain scales with how badly each configuration was starved is the strongest
available evidence that the diagnosis is the right one, because the budget figures predicted
that pattern before the experiment was run. The comparison between the two learning rates
then reverses.

\begin{table}[htp]
\centering
\begin{tabular}{lrr}
\toprule
{\bf Early-stopping rule} & {\bf Unreconciled} & {\bf With MinT} \\
\midrule
Original (validation block) & $+0.51$\% [$-0.16$, $+1.22$] & $+0.66$\% [$-0.12$, $+1.51$] \\
Aligned (training-window tail) & $-0.96$\% [$-1.39$, $-0.50$] & $-0.17$\% [$-0.84$, $+0.47$] \\
\bottomrule
\end{tabular}
\caption[Learning rate under each rule]{\textbf{$\eta = 0.005$ against $\eta = 0.03$ under
each early-stopping rule.} Negative favours the lower learning rate. Under the original rule
the sign is wrong; under the aligned rule the lower learning rate is significantly better
unreconciled and indistinguishable after MinT.}
\label{tab:lrflip}
\end{table}

The cross-validation result was correct and the delivery protocol was wrong. Under the
aligned rule the lower learning rate is better by 0.96\% on the base forecasts, close to
what cross-validation had predicted, and after MinT the two cannot be separated.

\subsection{What was adopted, and where the gain came from}

Delivery uses a learning rate of 0.03 rather than 0.005, together with the aligned
early-stopping rule. After reconciliation, which is what the system delivers, the two
learning rates are indistinguishable, while the lower one needs 18,556 boosting rounds and
110 minutes of training against 3,296 rounds and 20 minutes. An application needing the
unreconciled base forecasts, where the difference is significant, would choose the lower
rate.

\begin{table}[htp]
\centering
\begin{tabular}{lrrr}
\toprule
{\bf Version} & {\bf Unreconciled} & {\bf With MinT} & {\bf MinT MAE} \\
\midrule
Before tuning & $-32.9$\% & $-35.2$\% & 9,110 \\
Round-one winner, original rule & $-33.6$\% & $-36.6$\% & 8,906 \\
Round-two winner, original rule & $-33.2$\% & $-36.2$\% & 8,965 \\
Round-one winner, aligned rule (adopted) & $-34.2$\% & {\bf $-37.4$\%} & {\bf 8,797} \\
Round-two winner, aligned rule & $-34.9$\% & $-37.5$\% & 8,781 \\
\bottomrule
\end{tabular}
\caption[Tuning summary]{\textbf{The five model versions evaluated on the test period.}
Percentages are against \naiveA{}.}
\label{tab:tunesummary}
\end{table}

Against the configuration in place before any tuning, the adopted version improves pooled
MAE by 3.61\% with a paired interval of [$-4.69$, $-2.47$] after reconciliation, winning on
75 of the 107 origins. Of the 2.2~\pp{} of skill separating the two, roughly 1.4 came from
the round-one search, 0.8 from aligning the early-stopping rule, and 0.1 from the round-two
search, which the intervals cannot distinguish from zero.

The largest single gain in this group came from repairing a defect in the training code
rather than from choosing better hyperparameters, and the defect only became visible because
a configuration that had won on cross-validation lost on deployment. That is the same shape
as the result which opened the chapter, where a feature ranking that held on one validation
block reversed on the test period, and both belong to the discussion in
Chapter~\ref{ch:discussion}.
